\begin{table}[!ht]\centering
\caption{Post-2016 change in the automation-risk gradient of log hourly pay}\label{tab:pay}
\begin{threeparttable}\scriptsize\setlength{\tabcolsep}{2.5pt}
\begin{tabular}{lcccccc}\toprule
 & (1) & (2) & (3) & (4) & (5) & (6) \\
 & Mean pay & Balanced & 2014--19 & Weighted & Median pay & Excl.\ overtime \\ \midrule
Risk $\times$ post-2016 & 0.122 & 0.124 & 0.116 & 0.187 & 0.139 & 0.124 \\
 & (0.019) & (0.018) & (0.018) & (0.025) & (0.019) & (0.019) \\
95\% confidence interval & [0.085, 0.158] & [0.088, 0.161] & [0.082, 0.151] & [0.138, 0.236] & [0.101, 0.177] & [0.088, 0.161] \\ \addlinespace
Occupation-years & 2,546 & 2,506 & 2,184 & 2,329 & 2,496 & 2,544 \\
Occupations & 366 & 358 & 366 & 333 & 367 & 366 \\
Within $R^{2}$ & 0.029 & 0.032 & 0.035 & 0.172 & 0.046 & 0.031 \\
\bottomrule\end{tabular}
\begin{tablenotes}\footnotesize
\item \emph{Notes:} The dependent variable is the log of the stated pay measure and the specification is equation (1). All columns include occupation and release fixed effects. Standard errors clustered by occupation are reported in parentheses. Column 2 keeps the 358 occupations with a published mean in all seven releases. Column 4 weights each cell by the occupation's employee jobs in the 2015 release, on the 333 occupations with a published count; the unweighted estimate on that same sample is 0.134 (0.019). Column 5 uses the published median of Table 14.5a and column 6 the mean of Table 14.6a. All other columns weight each occupation-year equally.
\item \emph{Source:} ASHE Table 14, 2014--2020 releases, revised editions; ONS (2019) probability of automation.
\end{tablenotes}\end{threeparttable}\end{table}